% EDITA AQUÍ: datos de la tesis y color. Después, sustituye los capítulos de ejemplo.
% Edita cada capítulo en su archivo dentro de sections/.
\documentclass[11pt,a4paper,oneside]{report}
\usepackage[T1]{fontenc}
\usepackage[spanish,es-nodecimaldot,es-noquoting]{babel}
\usepackage{tgpagella}
\usepackage[inner=30mm,outer=25mm,top=27mm,bottom=27mm]{geometry}
\usepackage{microtype,amsmath,amssymb,graphicx,booktabs,tabularx}
\usepackage{xcolor,titlesec,fancyhdr,enumitem}
\definecolor{accent}{HTML}{245C61}
\definecolor{muted}{HTML}{59676B}
\usepackage[colorlinks=true,allcolors=accent]{hyperref}
\newcommand{\thesistitle}{Título de la tesis}
\newcommand{\thesissubtitle}{Un subtítulo descriptivo para tu investigación}
\newcommand{\authorname}{Tu nombre}
\newcommand{\university}{Tu universidad}
\newcommand{\department}{Departamento / Facultad}
\newcommand{\degree}{Trabajo de máster / Tesis doctoral}
\newcommand{\supervisor}{Nombre del director}
\hypersetup{pdftitle={\thesistitle},pdfauthor={\authorname}}
\linespread{1.12}
\setlength{\parindent}{0pt}\setlength{\parskip}{7pt}
\setlist[itemize]{leftmargin=*,itemsep=3pt,topsep=4pt}
\titleformat{\chapter}[display]{\bfseries\color{accent}}{\small\sffamily\MakeUppercase{\chaptertitlename}\ \thechapter}{8pt}{\Huge}
\titlespacing*{\chapter}{0pt}{0pt}{24pt}
\titleformat{\section}{\large\bfseries\color{accent}}{\thesection}{0.7em}{}
\pagestyle{fancy}\fancyhf{}
\fancyhead[L]{\small\color{muted}\nouppercase{\leftmark}}
\fancyfoot[L]{\small\color{muted}\authorname}
\fancyfoot[R]{\small\thepage}
\renewcommand{\headrulewidth}{0pt}\setlength{\headheight}{16pt}

\begin{document}
\hypersetup{pageanchor=false}
\begin{titlepage}
\raggedright
{\large\sffamily\color{accent}\university}\par
{\small\department}\par
\vspace{28mm}
{\small\sffamily\color{muted}\MakeUppercase{\degree}}\par
\vspace{8mm}
{\Huge\bfseries\thesistitle\par}
\vspace{6mm}
{\Large\color{muted}\thesissubtitle\par}
\vspace{12mm}{\color{accent}\hrule height 1.5pt}
\vspace{14mm}
{\large\authorname}\par
\vfill
\begin{tabular}{@{}ll@{}}
\textbf{Dirección} & \supervisor\\[5pt]
\textbf{Programa} & Tu programa\\[5pt]
\textbf{Entrega} & \today
\end{tabular}
\end{titlepage}
\pagenumbering{roman}
\hypersetup{pageanchor=true}

\chapter*{Resumen}
\addcontentsline{toc}{chapter}{Resumen}
Resume el problema, la laguna en el conocimiento y la pregunta principal.
Describe el método, el resultado más relevante y la aportación del trabajo.
Escribe esta sección al final para que refleje la tesis que has realizado.

Este ejemplo utiliza un estudio de procesos de investigación como hilo conductor.
Sustituye todos los ejemplos, incluida la tabla y la referencia, por tus propios
datos. Respeta el límite de palabras y los requisitos de tu institución.

\textbf{Palabras clave:} pregunta de investigación; metodología; evaluación; tu disciplina.
\vspace{12mm}

\section*{Agradecimientos}
Reconoce la dirección, la colaboración, la financiación y el apoyo personal.
Comprueba si tu institución exige una declaración de originalidad u otros preliminares.

\tableofcontents
\clearpage
\pagenumbering{arabic}

\chapter{Introducción y antecedentes}
\label{ch:introduction}
\section{Contexto y motivación}
Explica el problema en su contexto. Identifica a quién afecta, qué se puede hacer
actualmente y qué dificultad motiva la investigación. Apoya las afirmaciones
con fuentes adecuadas, por ejemplo \cite{example}.

Avanza del contexto general a una laguna concreta de conocimiento. Distingue
lo que demuestran estudios anteriores de lo que investigará tu tesis y explica
por qué esa laguna es relevante en la teoría o en la práctica.
\section{Pregunta de investigación y objetivos}
\textbf{Pregunta de investigación:} ¿Cómo afecta un procedimiento revisado al tiempo
de realización y a la precisión frente al proceso habitual?
\begin{itemize}
\item Establecer una referencia de comparación y criterios de evaluación explícitos.
\item Desarrollar o seleccionar un enfoque que aborde la laguna identificada.
\item Evaluar los resultados y discutir los límites de las pruebas.
\end{itemize}
\section{Trabajos relacionados}
Organiza la bibliografía por temas o enfoques. Compara sus supuestos, métodos
y hallazgos. Explica cómo tu estudio se deriva de esa comparación.
\section{Alcance y estructura de la tesis}
Delimita la población, el contexto, los materiales y los resultados del estudio.
El capítulo~\ref{ch:methods} describe el método; el capítulo~\ref{ch:findings} presenta
y discute los hallazgos. El apéndice~\ref{app:materials} contiene material complementario.

\chapter{Diseño de la investigación y métodos}
\label{ch:methods}
\section{Diseño y justificación}
Explica por qué el diseño responde a la pregunta de investigación. Identifica
la unidad de análisis, las condiciones de comparación y los controles.
Describe la selección de participantes y los criterios de inclusión y exclusión.
\section{Materiales y procedimiento}
Documenta los materiales, las herramientas y sus versiones. Describe cada fase
con suficiente detalle para reproducirla.
\begin{table}[htbp]
\centering
\caption{Protocolo de ejemplo; adáptalo a tu investigación.}
\begin{tabularx}{\linewidth}{@{}lX@{}}
\toprule
\textbf{Fase} & \textbf{Información que registrar}\\
\midrule
Preparación & Tarea, materiales y criterios de selección.\\
Recogida & Condiciones, observaciones y desviaciones del protocolo.\\
Análisis & Métricas, incertidumbre, exclusiones y código de análisis.\\
\bottomrule
\end{tabularx}
\end{table}
\section{Medidas y plan de análisis}
Define los resultados, las unidades y la comparación antes de presentar los datos.
Para el tiempo de realización, una posible medida es
\begin{equation}
\Delta_t = \frac{t_{\mathrm{habitual}}-t_{\mathrm{revisado}}}{t_{\mathrm{habitual}}}\times100\%.
\end{equation}
Justifica el análisis estadístico o cualitativo y el tratamiento de datos ausentes.
Distingue los análisis confirmatorios de los exploratorios.
\section{Ética y reproducibilidad}
Describe el consentimiento, la privacidad, la gestión de datos y las autorizaciones.
Indica cómo compartirás los materiales, datos y código, o justifica las restricciones.

\chapter{Hallazgos, discusión y conclusiones}
\label{ch:findings}
\section{Resultados}
Presenta las pruebas antes de interpretarlas. Describe la muestra, los resultados
principales y la incertidumbre. La tabla~\ref{tab:outcomes} contiene valores
ilustrativos; sustitúyelos por tus datos.
\begin{table}[htbp]
\centering
\caption{Comparación ilustrativa; no son resultados de investigación.}
\label{tab:outcomes}
\begin{tabularx}{\linewidth}{@{}Xrr@{}}
\toprule
\textbf{Medida} & \textbf{Habitual} & \textbf{Revisado}\\
\midrule
Tiempo de realización (minutos) & 24.0 & 18.0\\
Tasa de errores (\%) & 8.0 & 5.0\\
\bottomrule
\end{tabularx}
\end{table}

\section{Interpretación y aportación}
Responde a la pregunta de investigación. Relaciona las pruebas con los objetivos
del capítulo~\ref{ch:introduction} y con trabajos anteriores. Explica qué nuevo
conocimiento o capacidad aporta el estudio.
\section{Limitaciones y trabajo futuro}
Discute las amenazas a la validez, las limitaciones de muestreo y los supuestos
de medida. Separa las afirmaciones respaldadas por datos de las explicaciones
que requieren más investigación. Propón estudios de seguimiento concretos.

\section{Conclusión}
Resume la respuesta, la aportación y sus implicaciones. Esta sección debe poder
leerse por sí sola sin introducir pruebas ni afirmaciones nuevas.

\renewcommand{\bibname}{Referencias}
\cleardoublepage\phantomsection
\addcontentsline{toc}{chapter}{Referencias}

\begin{thebibliography}{9}
\bibitem{example} A. Autor. \textit{Título de una obra relevante}. Editorial, 2025.
% Sustituye esta referencia de ejemplo por una fuente real. Añade entradas con \bibitem.
\end{thebibliography}

\appendix
\chapter{Material complementario}
\label{app:materials}
\section{Lista de reproducibilidad}
Incluye instrumentos, resultados ampliados, demostraciones o detalles de
implementación que apoyen el argumento principal sin interrumpirlo.
\begin{itemize}
\item Materiales e instrucciones para repetir el estudio.
\item Diccionario de datos, unidades y decisiones de anonimización.
\item Versiones de software y pasos del análisis.
\item Tablas u observaciones adicionales citadas en el texto.
\end{itemize}
\section{Ejemplo de registro de datos}
\begin{tabularx}{\linewidth}{@{}lX@{}}
\toprule
\textbf{Campo} & \textbf{Significado}\\
\midrule
ID de observación & Identificador anónimo de cada tarea registrada.\\
Condición & Procedimiento habitual o revisado.\\
Duración & Tiempo transcurrido, en minutos.\\
\bottomrule
\end{tabularx}

\end{document}
